\begin{SCfigure}[][b]
	\includegraphics[width=0.5\linewidth]{figures/logical_error_probability.pdf}
	\caption{The logical error rate per kept shot of the double-check circuit,
	plotted against noise level $p$.
	Dashed (solid) lines correspond to distance-3 (-5) cultivation,
	and blue (orange) data correspond to cultivating S (T) states.
	These curves are analytic and correct through fifth order in $p$.
	From top to bottom, the gradient of the lines as $p \to 0$ are: 2, 3, 3, 5.
	This figure and \cref{fig:lep_against_fault_count} are derived from the same enumeration data.}
	\label{fig:ler_against_per}
\end{SCfigure}

\section{Fault Enumeration of Magic State Cultivation}\label{sec:thorough_analysis}
We resume our focus on MSC,
combining \cref{lem:insert_projector,sec:clifford_error_analysis}
to find fault configurations in the final double check
with some probability of being undetected;
that is,
not flipping any detectors during and after the double check.
We use the same uniform depolarising noise model as
in the original MSC paper~\cite{Gidney2024},
which is summarised in \cref{sec:noise_model}.
We enumerate all such configurations of \num{\le 5} faults
for both distances 3 and 5;
the latter took 15 seconds on a 2021 Apple MacBook Pro M1.
Compared to Monte Carlo sampling,
enumeration gives us more insight
and allows us to analytically reconstruct
an approximation of the logical error rate not at one noise level,
but as a function of the noise level.
We need only compute this function once
before we can instantaneously probe arbitrarily low logical error rates
and generate plots such as \cref{fig:ler_against_per}.
We note this approach is similar to SyQMA~\cite{Umbrarescu2026},
which can compute the exact symbolic logical error rate
(and thus fault distance) of non-Clifford circuits,
but does so by simulating the noisy state rather than enumerating fault configurations.

We say an undetected fault configuration is
benign (malignant)
if it leads to a logical fidelity of 1 (0).
For MSC-3,
there are 4 malignant 2-fault configurations;
these are explicitly shown in \cref{fig:all_2_fault_configurations}.
\begin{SCfigure}
	\begingroup

\definecolor{mscCircle}{HTML}{0072B2}
\definecolor{mscSquare}{HTML}{E69F00}
\definecolor{mscDiamond}{HTML}{009E73}
\definecolor{mscTriangle}{HTML}{CC79A7}

\tikzset{
	msc circle wire/.style={msc wire above=mscCircle},
	msc square wire/.style={msc wire above=mscSquare},
	msc diamond wire/.style={msc wire below=mscDiamond},
	msc triangle wire/.style={msc wire above=mscTriangle},
	msc circle square wire/.style={
		msc wire above=mscCircle,
		msc wire below=mscSquare,
	},
	msc square diamond wire/.style={
		msc wire above=mscSquare,
		msc wire below=mscDiamond,
	},
	msc circle control/.style={msc control left=mscCircle},
	msc square control/.style={msc control left=mscSquare},
	msc diamond control/.style={msc control left=mscDiamond},
	msc triangle control/.style={msc control left=mscTriangle},
	msc circle square control/.style={
		msc control left=mscCircle,
		msc control right=mscSquare,
	},
	msc square diamond control/.style={
		msc control left=mscSquare,
		msc control right=mscDiamond,
	},
	msc circle event/.style={
		msc event,
		shape=yquant-circle,
		fill=mscCircle!25,
		draw=mscCircle!85!black,
	},
	msc square event/.style={
		msc event,
		fill=mscSquare!25,
		draw=mscSquare!85!black,
	},
	msc diamond event/.style={
		msc event,
		shape=diamond,
		fill=mscDiamond!25,
		draw=mscDiamond!85!black,
		inner timesep=0.3mm,
		inner spacesep=0.3mm,
	},
	msc triangle event/.style={
		msc event,
		shape=regular polygon,
		regular polygon sides=3,
		fill=mscTriangle!25,
		draw=mscTriangle!85!black,
		inner timesep=0.1mm,
		inner spacesep=0.1mm,
	},
}

\yquantset{
	msc circle cnot/.style={
		every control line/.append style={msc circle control},
	},
	msc square cnot/.style={
		every control line/.append style={msc square control},
	},
	msc diamond cnot/.style={
		every control line/.append style={msc diamond control},
	},
	msc triangle cnot/.style={
		every control line/.append style={msc triangle control},
	},
	msc circle square cnot/.style={
		every control line/.append style={msc circle square control},
	},
	msc square diamond cnot/.style={
		every control line/.append style={msc square diamond control},
	},
}

\begin{tikzpicture}[trim left=10.5pt]
\begin{yquant}

qubit {0} q0;
qubit {a} qa;
qubit {b} qb;
qubit {1} q1;
qubit {c} qc;
qubit {2} q2;
qubit {d} qd;
qubit {3} q3;
qubit {4} q4;
qubit {e} qe;
qubit {5} q5;
qubit {6} q6;
qubit {f} qf;

discard qa,qb,qc,qd,qe,qf;
hspace {1mm} q0,qa,qb,q1,qc,q2,qd,q3,q4,qe,q5,q6,qf;
box {T} q1,q4,q5;
box {T$^\dag$} q0,q2,q3,q6;
init {$\ket{+}$} qa,qb,qc,qd,qe,qf;

cnot q0|qa;
cnot q4|qe;
cnot q3|qd;
cnot q2|qc;
cnot q1|qb;
cnot q6|qf;

cnot qa|q1;
cnot qd|q2;
cnot qf|qe;

[msc circle event]
box {X} q1;
[msc square event]
box {X} q2;
[msc diamond event]
box {X} qe;
setstyle {msc circle wire} q1;
setstyle {msc square wire} q2;
setstyle {msc diamond wire} qe;
align q0,qa,qb,q1,qc,q2,qd,q3,q4,qe,q5,q6,qf;

[msc square cnot]
cnot q1|q2;
setstyle {msc circle square wire} q1;
[msc diamond cnot]
cnot q5|qe;
setstyle {msc diamond wire} q5;

[msc square cnot]
cnot qe|q2;
setstyle {msc square diamond wire} qe;

dmeter {X} q2;
discard q2;
init {$\ket{+}$} q2;
setstyle {} q2;

cnot qe|q2;

cnot q1|q2;
[msc square diamond cnot]
cnot q5|qe;
setstyle {msc square wire} q5;

[msc triangle event]
box {X} q2;
setstyle {msc triangle wire} q2;
align q0,qa,qb,q1,qc,q2,qd,q3,q4,qe,q5,q6,qf;

[msc circle square cnot]
cnot qa|q1;
setstyle {msc circle square wire} qa;
[msc triangle cnot]
cnot qd|q2;
setstyle {msc triangle wire} qd;
[msc square diamond cnot]
cnot qf|qe;
setstyle {msc square diamond wire} qf;

[msc circle square cnot]
cnot q0|qa;
setstyle {msc circle square wire} q0;
[msc square diamond cnot]
cnot q4|qe;
setstyle {msc square diamond wire} q4;
[msc triangle cnot]
cnot q3|qd;
setstyle {msc triangle wire} q3;
cnot q2|qc;
cnot q1|qb;
[msc square diamond cnot]
cnot q6|qf;
setstyle {msc square diamond wire} q6;

[msc circle event]
box {Y} q3;
[msc square event]
box {Y} q4;
[msc triangle event]
box {Y} q4;
[msc diamond event]
box {Y} q0;

dmeter {X} qa,qb,qc,qd,qe,qf;
discard qa,qb,qc,qd,qe,qf;
align -;
box {T$^\dag$} q1,q4,q5;
box {T} q0,q2,q3,q6;
setstyle {} q0,q1,q2,q3,q4,q5,q6;

hspace {1mm} q0,q1,q2,q3,q4,q5,q6;

\end{yquant}
\end{tikzpicture}

\endgroup

	\caption{All four 2-fault configurations
	for the distance-3 double-check circuit (\cref{fig:d3a6})
	that can lead to logical error.
	Each configuration is shown as a pair of error events
	of a given shape and colour
	e.g.\ the blue circled X and Y realise
	the configuration in \cref{fig:2_fault_configuration}.
	Note each configuration can be realised by
	multiple equivalent pairs of error events;
	here we choose to show the earliest possible events.
	We also highlight how each X event propagates until the layer of T-gates.}
	\label{fig:all_2_fault_configurations}
\end{SCfigure}
Note these configurations degrade the fault distance of MSC-3 only:
they are caught by the subsequent distance-5 double check in MSC-5.
For MSC-5,
there are 3 (601) malignant 3- (4-) fault configurations
in its final, distance-5 double check;
each configuration contains at least one X hook error.
\Cref{fig:fcc_against_fault_count} shows
the complete data on configuration counts,
from which \cref{fig:lep_against_fault_count} is derived.
In \cref{sec:logical_error_rate_per_kept_shot} we describe how to
analytically construct the logical error rate per kept shot
from these fault enumerations,
the result of which is shown in \cref{fig:ler_against_per}.
The T--S logical error rate discrepancy
observed for both code distances in this subfigure
can be explained by \cref{fig:lep_against_fault_count}.
The first reason is the difference in fault distance
(indicated by the first nonzero point for each line);
this dictates the asymptotic scaling
of the logical error rate as the noise level $p \to 0$.
The second reason is that T-state cultivation has
a higher fraction of malignant configurations
compared to S-state cultivation;
this manifests as a higher constant factor
multiplying the logical error rate.

This degraded fault distance is experienced also by
the distance-3 double check used in
the experimental realisation of MSC~\cite{Rosenfeld2025}
(Crumble link \href{https://algassert.com/crumble\#circuit=Q(0,1)0;Q(0,2)1;Q(0.5,0.5)2;Q(0.5,1.5)3;Q(1,0)4;Q(1,1)5;Q(1,2)6;Q(1.5,0.5)7;Q(1.5,1.5)8;Q(2,1)9;Q(2,2)10;POLYGON(0,0,1,0.25)0_4_9_5;POLYGON(0,1,0,0.25)0_5_6_1;POLYGON(1,0,0,0.25)5_9_10_6;TICK;MPP_Y0*Y4*Y5*Y9*Y10*Y6*Y1;TICK;MPP_X5*X9*X10*X6;TICK;MPP_X0*X5*X6*X1;TICK;MPP_X4*X9*X5*X0;TICK;MPP_Z5*Z9*Z10*Z6;TICK;MPP_Z0*Z5*Z6*Z1;TICK;MPP_Z4*Z9*Z5*Z0;TICK;RX_8_2_7_3;TICK;S_0_9_6;S_DAG_4_5_1_10;TICK;CX_3_0_7_4_2_5_8_10;TICK;CX_5_7_8_9_3_1;TICK;CX_5_3_8_6;TICK;CX_8_5;TICK;MX_8;DT(0,1,0)rec[-1]_rec[-8];TICK;RX_8;TICK;CX_8_5;TICK;CX_5_3_8_6;TICK;CX_5_7_8_9_3_1;TICK;CX_3_0_7_4_2_5_8_10;TICK;S_4_5_1_10;S_DAG_0_9_6;TICK;MX_8_2_7_3;DT(1.5,1.5,1)rec[-4]_rec[-5];DT(0.5,0.5,1)rec[-3];DT(1.5,0.5,1)rec[-2];DT(0.5,1.5,1)rec[-1];TICK;MPP_X5*X9*X10*X6;DT(1,1,2)rec[-1]_rec[-12];TICK;MPP_X0*X5*X6*X1;DT(0,1,3)rec[-1]_rec[-12];TICK;MPP_X4*X9*X5*X0;DT(1,0,4)rec[-1]_rec[-12];TICK;MPP_Z5*Z9*Z10*Z6;DT(1,1,5)rec[-1]_rec[-12];TICK;MPP_Z0*Z5*Z6*Z1;DT(0,1,6)rec[-1]_rec[-12];TICK;MPP_Z4*Z9*Z5*Z0;DT(1,0,7)rec[-1]_rec[-12];TICK;MPP_Y0*Y4*Y5*Y9*Y10*Y6*Y1;OI(0)rec[-1]_rec[-8]_rec[-9]_rec[-11]}{here}):
applying the same analysis to this slightly smaller circuit,
we find it also has fault distance 2.
